\honest{Evolutions Are Stupid (But Work Anyway)}{Eliezer Yudkowsky}{2007}

Yesterday I wrote that ``Science has a very exact idea of the capabilities of evolution.'' Today I show what I meant. I say ``evolutions,'' plural, because each population evolves on its own. None of this is creationism; it is ``standard Evolutionary Biology 201,'' and the limits are needed to make sense of observed biology. Human intelligence is too complicated to measure its efficiency, but natural selection, the first optimization process, is simpler, slower and less efficient.

Evolutions are simple enough that ``we can calculate exactly how stupid they are.'' A gene with a 3 per cent advantage takes about 768 generations to spread through a population of 100,000, and 875 through 500,000. The formula is $2\ln(N)/s$, where $N$ is the population and $1+s$ the relative fitness. In a hunter-gatherer population of a million, a 1 per cent advantage takes 2,763 generations. The chance that it spreads at all is about twice its advantage, 6 per cent, whatever the population size. The same mutation may recur, but in a population of a million you may wait a hundred generations for another try, again at 6 per cent. Still, in the long run an evolution has a good chance of getting there. \nb{The arithmetic is right, and the figures agree to within about 5 per cent with the standard results of Kimura and Ohta and of Haldane. Both are approximations for an idealized population of constant size, good for orders of magnitude rather than exact.} Genes have no investors or imitators; they spread only by their bearers having more children. It is as if Henry Ford had to make one car, sell it, buy parts for 1.01 more cars, and repeat until he had a million.

Complex adaptations take far longer. A fur coat helps only if the environment reliably brings cold, and genes are part of each other's environment. If B gives a 5 per cent advantage in the presence of A and none otherwise, then while A is at 1 per cent B's average advantage is 0.05 per cent and its chance of fixation 0.1 per cent. So ``Only then,'' once A has spread, can B begin to rise. \nb{This assumes B appears as a new mutation after A has spread. A gene already present at low frequency starts rising as soon as A is common, so the steps can overlap and adaptation is faster.} A complex adaptation takes a thousand generations for A, another thousand for B, another for A*, and millions of years in all. Then ``other evolutions don't imitate it'': a snake's new venom does not help the fox. \nb{True of snakes and foxes. Bacteria pass genes between species, and antibiotic resistance spreads this way.}

A human programmer can design a hundred interdependent parts in an afternoon. How is that possible? I do not fully know, and I suspect science does not. Human brains are much more complicated than evolutions. But humans plan ahead, change several parts together, learn from single cases, think abstractly about problem spots and choose which changes to try instead of waiting for a cosmic ray. By natural selection's standards this is magic. I quote the biologist Cynthia Kenyon, at a dinner, saying that a graduate student can do in an hour what evolution could not do in a billion years. Evolutions have invented the rotating wheel three times. I give no source for the count. Then the programmer posts the code to the Internet.

Some of biology is impressive even beside our best technology: we cannot build a self-replicating machine the size of a butterfly. But we only began accumulating knowledge about four hundred years ago, and we already have wheels, steel, rockets, transistors and nuclear power. The balance tips further every decade. So, once again, for a human to learn design from natural selection would be like a modern bacterium imitating the first replicator, which would be eaten at once in today's ecology. So would a human planner who made random changes to a strategy and waited 768 rounds of testing to adopt a 3 per cent improvement. Don't praise evolutions ``one millimeter more than they deserve.''

I promise ``More exciting mathematical bounds on evolution'' for tomorrow. \nb{That post, ``Natural Selection's Speed Limit and Complexity Bound,'' is not in the book. It now carries Yudkowsky's warning that a simulation written to check it ``did not return the expected results.''}
